\documentclass[12pt]{article}

\usepackage[a4paper,margin=1in]{geometry}
\usepackage{amsmath,amssymb}
\usepackage{newtxtext,newtxmath}
\usepackage{fancyhdr}
\usepackage{listings}
\usepackage{graphicx}
\usepackage{float}
\usepackage{hyperref}
\usepackage{booktabs}
\usepackage{xcolor}

\setlength{\parindent}{0pt}
\setlength{\parskip}{5pt}

% Keep entire document at 12 pt
\AtBeginDocument{
    \fontsize{12pt}{14pt}\selectfont
}

% ------------------------------------------------
% Header Configuration
% ------------------------------------------------
\pagestyle{fancy}
\fancyhf{}
\fancyhead[L]{Name: \underline{Kishan Sahu}}
\fancyhead[R]{Enrollment No.: \underline{P26DS017}}
\renewcommand{\headrulewidth}{0.4pt}
\setlength{\headheight}{15pt}

% ------------------------------------------------
% Code Formatting Configuration
% ------------------------------------------------
\lstset{
    basicstyle=\ttfamily\fontsize{8.5pt}{10.5pt}\selectfont,
    breaklines=true,
    frame=single,
    columns=fullflexible,
    keepspaces=true,
    showstringspaces=false,
    keywordstyle=\color{blue}\bfseries,
    commentstyle=\color{gray}\itshape,
    stringstyle=\color{teal}
}

\begin{document}

% ==============================================================================
% TITLE BANNER
% ==============================================================================
\begin{center}
\textbf{\Large SARDAR VALLABHBHAI NATIONAL INSTITUTE OF TECHNOLOGY, SURAT}\\[4pt]
\textbf{\large Department of Computer Science \& Engineering}\\[3pt]
\textbf{\large CS103 / CSDS103: Design and Analysis of Algorithms}\\[3pt]
\textbf{\large M.Tech. (DS) I Year -- Semester I}\\[4pt]
\textbf{\Large Lab Assignment 11: Graph Algorithms -- Minimum Arborescence, Shortest Paths, Arbitrage Detection, DAG Scheduling \& BFS Navigation}
\end{center}

\vspace{0.2cm}
\hrule
\vspace{0.4cm}

\textbf{Overview \& Scope}:\\
This laboratory report presents an end-to-end mathematical, algorithmic, and empirical investigation of five canonical graph optimization problems:
\begin{enumerate}
    \item \textbf{Assignment 1}: Design a Minimum-Cost One-Way Delivery Network Using Minimum Arborescence (Chu-Liu / Edmonds' Algorithm in $\mathcal{O}(V \cdot E)$ Time).
    \item \textbf{Assignment 2}: Emergency Route Finder Using Dijkstra's Algorithm with Min-Heap Priority Queue in $\mathcal{O}((V + E) \log V)$ Time.
    \item \textbf{Assignment 3}: Currency Exchange Arbitrage Detection Using the Bellman--Ford Algorithm via Logarithmic Weight Transformation ($w(u, v) = -\ln R(u, v)$).
    \item \textbf{Assignment 4}: Project Scheduling Using a Directed Acyclic Graph (DAG) via Kahn's Topological Sorting and Critical Path Analysis.
    \item \textbf{Assignment 5}: Campus Walkway Navigation Using Breadth-First Search (BFS) for Fewest Walkway Segments on an Unweighted Graph.
\end{enumerate}

\vspace{0.3cm}

% ==============================================================================
% ASSIGNMENT 1: MINIMUM ARBORESCENCE
% ==============================================================================
\section*{Assignment 1: Minimum-Cost One-Way Delivery Network Using Minimum Arborescence}

\subsection*{1.1 Problem Statement and Mathematical Formulation}
\textbf{Problem}: A logistics enterprise operates a central distribution warehouse and multiple regional distribution centers. Candidate delivery routes are strictly one-way (directed) and carry distinct transportation costs. We are required to design a minimum-cost directed distribution network such that:
\begin{enumerate}
    \item Every regional distribution center is reachable from the central warehouse via directed paths.
    \item Each non-root distribution center has in-degree exactly $1$ in the chosen subgraph.
\end{enumerate}

\textbf{Mathematical Formulation}:
Let $G = (V, E)$ be a directed weighted graph where $V$ represents the set of logistics nodes ($|V| = n$), $r \in V$ is the designated root (central warehouse), and $w: E \to \mathbb{R}^+$ assigns non-negative costs to one-way delivery routes. A subgraph $T = (V, E_T)$ is called a \textit{directed spanning tree} (or \textit{arborescence}) rooted at $r$ if:
\begin{itemize}
    \item $T$ contains no directed cycles.
    \item The in-degree of the root is zero: $\text{deg}^-_T(r) = 0$.
    \item Every other vertex $v \in V \setminus \{r\}$ has in-degree exactly one: $\text{deg}^-_T(v) = 1$.
    \item For every $v \in V$, there exists a unique directed path from $r$ to $v$.
\end{itemize}
The objective is to minimize total network infrastructure cost:
\[
\min_{E_T \subseteq E} \sum_{e \in E_T} w(e)
\]

\textbf{Algorithmic Principle (Chu-Liu / Edmonds' Algorithm)}:
Unlike undirected Minimum Spanning Trees (which can be solved by Kruskal's or Prim's greedy edge selection), a purely greedy incoming-edge selection on directed graphs can form isolated directed cycles. The Chu-Liu / Edmonds algorithm resolves this via recursive cycle contraction:
\begin{enumerate}
    \item \textbf{Greedy Choice}: For each non-root vertex $v \in V \setminus \{r\}$, select the incoming edge $(u, v)$ minimizing $w(u, v)$. Let this edge be denoted by $\text{in}(v)$.
    \item \textbf{Cycle Check}: If the selected edges form no directed cycle, the edge set $\{\text{in}(v) \mid v \in V \setminus \{r\}\}$ is provably the optimal minimum spanning arborescence.
    \item \textbf{Cycle Contraction}: If a directed cycle $C$ is detected:
    \begin{itemize}
        \item Contract all vertices of $C$ into a single composite pseudo-node $v_C$.
        \item For every edge entering the cycle from outside ($u \notin C, v \in C$), adjust its effective weight:
        \[
        w'(u, v_C) = w(u, v) - w(\text{in}(v))
        \]
        This represents the marginal cost of breaking cycle $C$ at vertex $v$ to admit the external incoming path.
        \item Solve recursively on the contracted graph $G'$.
        \item Expand the contracted solution: if the optimal solution on $G'$ includes an edge entering $v \in C$, include that edge, break cycle $C$ by omitting $\text{in}(v)$, and retain all remaining $|C| - 1$ cycle edges.
    \end{itemize}
\end{enumerate}

\subsection*{1.2 Python Implementation (\texttt{min\_arborescence.py})}
\begin{lstlisting}[language=Python]
from typing import List, Tuple, Dict, Optional, Set

def find_minimum_arborescence(
    nodes: List[str],
    root: str,
    edges: List[Tuple[str, str, float]]
) -> Tuple[Optional[float], List[Tuple[str, str, float]]]:
    node_to_idx = {name: i for i, name in enumerate(nodes)}
    idx_to_node = {i: name for i, name in enumerate(nodes)}
    r_idx = node_to_idx[root]
    raw_edges = [(node_to_idx[u], node_to_idx[v], cost, i) for i, (u, v, cost) in enumerate(edges)]

    def edmonds(n: int, r: int, edge_list: List[Tuple[int, int, float, int]]):
        # 1. Select cheapest incoming edge for each non-root node
        min_in: Dict[int, Tuple[float, int, int]] = {}
        for u, v, w, eid in edge_list:
            if u != v and v != r:
                if v not in min_in or w < min_in[v][0]:
                    min_in[v] = (w, u, eid)

        for v in range(n):
            if v != r and v not in min_in:
                return None, []  # Node unreachable from root

        # 2. Cycle detection in chosen minimum incoming edges
        comp = [-1] * n
        vis = [-1] * n
        num_cycles = 0
        cycle_edges = {}

        for i in range(n):
            if i == r:
                continue
            curr = i
            path = []
            while curr != r and curr != -1 and vis[curr] == -1:
                vis[curr] = i
                path.append(curr)
                curr = min_in[curr][1] if curr in min_in else -1

            if curr != -1 and curr != r and vis[curr] == i:
                cycle = []
                c = path.pop()
                cycle.append(c)
                while c != curr:
                    c = path.pop()
                    cycle.append(c)
                for node in cycle:
                    comp[node] = num_cycles
                    cycle_edges[node] = min_in[node]
                num_cycles += 1

        # Base Case: acyclic selection -> minimum arborescence found!
        if num_cycles == 0:
            sel_eids = [min_in[v][2] for v in range(n) if v != r]
            cost = sum(min_in[v][0] for v in range(n) if v != r)
            return cost, sel_eids

        # 3. Contract cycles into pseudo-vertices
        comp_count = num_cycles
        for v in range(n):
            if comp[v] == -1:
                comp[v] = comp_count
                comp_count += 1

        contracted_edges = []
        for u, v, w, eid in edge_list:
            cu, cv = comp[u], comp[v]
            if cu != cv:
                if cv < num_cycles:
                    contracted_edges.append((cu, cv, w - min_in[v][0], eid))
                else:
                    contracted_edges.append((cu, cv, w, eid))

        cost, rec_eids = edmonds(comp_count, comp[r], contracted_edges)
        if cost is None:
            return None, []

        # 4. Expand contracted cycle edges
        cycle_incoming = {}
        for eid in rec_eids:
            for u, v, w, orig_eid in edge_list:
                if orig_eid == eid:
                    cv = comp[v]
                    if cv < num_cycles:
                        cycle_incoming[cv] = v
                    break

        final_eids = set(rec_eids)
        for node, (cw, cu, ceid) in cycle_edges.items():
            cid = comp[node]
            if cid not in cycle_incoming or node != cycle_incoming[cid]:
                final_eids.add(ceid)

        return cost, list(final_eids)

    cost, sel_ids = edmonds(len(nodes), r_idx, raw_edges)
    if cost is None:
        return None, []

    selected_edges = [edges[i] for i in sel_ids]
    total_cost = sum(cost for u, v, cost in selected_edges)
    return total_cost, selected_edges
\end{lstlisting}

\subsection*{1.3 Theoretical Complexity Analysis}
\begin{itemize}
    \item \textbf{Cycle Detection \& Contraction Time}: Detecting cycles traverses $n$ nodes and takes $\mathcal{O}(V)$ time. Rewriting edges during contraction takes $\mathcal{O}(E)$ time.
    \item \textbf{Recursion Depth}: Each contraction strictly reduces the vertex count by at least $1$. Thus, there are at most $\mathcal{O}(V)$ recursive contraction phases.
    \item \textbf{Total Time Complexity}: $\mathcal{O}(V \cdot E)$. On sparse graphs where $E = \mathcal{O}(V)$, this runs in $\mathcal{O}(V^2)$ time.
    \item \textbf{Auxiliary Space Complexity}: $\mathcal{O}(V + E)$ to maintain component arrays, tracking structures, and edge identifiers.
\end{itemize}

\subsection*{1.4 Execution Output \& Verification}
Testing on a five-node delivery network with Central Warehouse and distribution centers:
\begin{lstlisting}
============================================================
ASSIGNMENT 1: MINIMUM ARBORESCENCE DELIVERY NETWORK
============================================================
Central Warehouse (Root): Central Warehouse
Regional Distribution Centers: ['DC North', 'DC South', 'DC East', 'DC West']

Minimum Total Network Delivery Cost: 30

Selected One-Way Delivery Routes (Arborescence Edges):
  Central Warehouse    ---> DC North        | Route Cost: 10
  DC East              ---> DC West         | Route Cost: 7
  DC North             ---> DC South        | Route Cost: 5
  DC South             ---> DC East         | Route Cost: 8
Verification Passed: All centers reachable from root with in-degree strictly 1!
\end{lstlisting}

\newpage

% ==============================================================================
% ASSIGNMENT 2: EMERGENCY ROUTE FINDER (DIJKSTRA)
% ==============================================================================
\section*{Assignment 2: Emergency Route Finder Using Dijkstra's Algorithm}

\subsection*{2.1 Problem Statement and Algorithmic Principle}
\textbf{Problem}: An emergency ambulance service stationed at a municipal depot needs to reach a critical care hospital in minimum travel time. City intersections are modeled as vertices, and connecting roadways are represented as directed edges weighted by positive expected travel times (accounting for traffic and transit delays).

\textbf{Mathematical Formulation}:
Let $G = (V, E)$ be a directed graph with edge travel-time weights $w: E \to \mathbb{R}^+_{\ge 0}$. Given origin $s$ (depot) and destination $t$ (hospital), find a directed simple path $P = (s = v_0, v_1, \dots, v_k = t)$ that minimizes:
\[
\text{Cost}(P) = \sum_{i=0}^{k-1} w(v_i, v_{i+1})
\]

\textbf{Dijkstra's Algorithm with Min-Heap Priority Queue}:
\begin{enumerate}
    \item Maintain a tentative distance array $\text{dist}[v]$, initialized to $\text{dist}[s] = 0$ and $\text{dist}[v] = \infty$ for all $v \neq s$.
    \item Insert $(0, s)$ into a min-priority queue (min-heap).
    \item While the heap is non-empty, extract $(d, u)$ with minimum travel time:
    \begin{itemize}
        \item If $d > \text{dist}[u]$, discard as an outdated entry.
        \item If $u = t$, the shortest path to the hospital is finalized.
        \item For each outgoing edge $(u, v)$ with weight $w(u, v)$, perform edge relaxation:
        \[
        \text{if } \text{dist}[u] + w(u, v) < \text{dist}[v] \implies \text{dist}[v] = \text{dist}[u] + w(u, v), \quad \text{parent}[v] = u
        \]
        and push $(\text{dist}[v], v)$ into the heap.
    \end{itemize}
\end{enumerate}

\textbf{Proof of Optimality}:
By induction on the set of settled vertices $S$. Since all edge weights are non-negative ($w(u, v) \ge 0$), any alternate path from $s$ to an unvisited node $u$ must pass through some boundary vertex $y \notin S$ with $\text{dist}[y] \ge \text{dist}[u]$. Because edge weights cannot be negative, no future edge additions can reduce $\text{dist}[u]$. Hence, $\text{dist}[u]$ is exact upon extraction.

\subsection*{2.2 Python Implementation (\texttt{emergency\_route.py})}
\begin{lstlisting}[language=Python]
import heapq
from typing import Dict, List, Tuple, Optional

def find_emergency_route(
    graph: Dict[str, List[Tuple[str, float]]],
    start: str,
    destination: str
) -> Tuple[Optional[float], List[str]]:
    pq: List[Tuple[float, str, List[str]]] = [(0.0, start, [start])]
    visited: Dict[str, float] = {}

    while pq:
        curr_time, u, path = heapq.heappop(pq)

        if u in visited and visited[u] <= curr_time:
            continue
        visited[u] = curr_time

        if u == destination:
            return curr_time, path

        for v, travel_time in graph.get(u, []):
            if v not in visited or curr_time + travel_time < visited[v]:
                heapq.heappush(pq, (curr_time + travel_time, v, path + [v]))

    return None, []
\end{lstlisting}

\subsection*{2.3 Theoretical Complexity Analysis}
\begin{itemize}
    \item \textbf{Heap Insertions and Extractions}: At most $|V|$ vertices are extracted, each requiring $\mathcal{O}(\log V)$ time. At most $|E|$ edge relaxations push new entries into the heap, each costing $\mathcal{O}(\log V)$.
    \item \textbf{Total Time Complexity}: $\mathcal{O}((V + E) \log V)$.
    \item \textbf{Auxiliary Space Complexity}: $\mathcal{O}(V)$ for the priority queue and distance/predecessor records.
\end{itemize}

\subsection*{2.4 Execution Output \& Verification}
Testing on the municipal road network:
\begin{lstlisting}
============================================================
ASSIGNMENT 2: EMERGENCY AMBULANCE ROUTE FINDER (DIJKSTRA)
============================================================
Origin:      Ambulance Station
Destination: City Hospital

Minimum Travel Time: 10.0 minutes
Optimal Emergency Route:
  Ambulance Station -> Intersection B -> Intersection C -> Intersection D -> City Hospital
Verification Passed: Path is strictly optimal and contains no backward loops.
\end{lstlisting}

\newpage

% ==============================================================================
% ASSIGNMENT 3: CURRENCY ARBITRAGE (BELLMAN-FORD)
% ==============================================================================
\section*{Assignment 3: Currency Exchange Arbitrage Detection Using Bellman--Ford}

\subsection*{3.1 Problem Statement and Algorithmic Principle}
\textbf{Problem}: A quantitative trading system processes a foreign exchange matrix of bilateral exchange rates. An arbitrage opportunity exists if a trader can convert an initial amount of currency through a sequence of foreign currencies and return to the home currency with strictly more capital than started with.

\textbf{Mathematical Formulation \& Logarithmic Duality}:
Let $C_1, C_2, \dots, C_n$ denote currencies. An exchange rate $R(u, v) > 0$ means $1$ unit of currency $u$ converts to $R(u, v)$ units of currency $v$. A trading cycle $c_1 \to c_2 \to \dots \to c_k \to c_1$ yields risk-free profit if and only if the cumulative product of exchange rates exceeds unity:
\[
\prod_{i=1}^{k} R(c_i, c_{i+1}) > 1.0 \quad (\text{where } c_{k+1} = c_1)
\]
Directly optimizing this product is non-linear and multiplicative. Taking the natural logarithm transforms multiplication into addition:
\[
\ln\left(\prod_{i=1}^{k} R(c_i, c_{i+1})\right) > 0 \iff \sum_{i=1}^{k} \ln R(c_i, c_{i+1}) > 0
\]
Negating both sides reverses the inequality:
\[
\sum_{i=1}^{k} \left(-\ln R(c_i, c_{i+1})\right) < 0
\]
Thus, by defining directed edge weights as:
\[
w(u, v) = -\ln R(u, v)
\]
the financial arbitrage detection problem reduces \textit{isomorphically} to detecting a \textbf{negative-weight cycle} in a directed graph!

\textbf{Bellman--Ford Negative Cycle Detection}:
\begin{enumerate}
    \item Initialize $\text{dist}[v] = 0.0$ for all currencies $v \in V$ (equivalent to introducing an artificial super-source with $0$-weight edges to every currency, ensuring all components are inspected).
    \item Relax all edges $|V| - 1$ times:
    \[
    \text{if } \text{dist}[u] + w(u, v) < \text{dist}[v] \implies \text{dist}[v] = \text{dist}[u] + w(u, v), \quad \text{parent}[v] = u
    \]
    \item Run an extra $|V|$-th relaxation pass: if any edge $(u, v)$ satisfies $\text{dist}[u] + w(u, v) < \text{dist}[v]$, a negative cycle exists.
    \item Backtrack predecessor pointers $|V|$ times starting from $v$ to guarantee entering a vertex inside the cycle, and traverse backwards until the cycle closes.
\end{enumerate}

\subsection*{3.2 Python Implementation (\texttt{currency\_arbitrage.py})}
\begin{lstlisting}[language=Python]
import math
from typing import List, Tuple, Dict, Optional

def detect_currency_arbitrage(
    currencies: List[str],
    rates: List[Tuple[str, str, float]]
) -> Tuple[Optional[List[str]], float]:
    edges = []
    rate_map = {}
    for u, v, r in rates:
        w = -math.log(r)
        edges.append((u, v, w))
        rate_map[(u, v)] = r

    n = len(currencies)
    dist = {c: 0.0 for c in currencies}
    parent = {c: None for c in currencies}

    for _ in range(n - 1):
        for u, v, w in edges:
            if dist[u] + w < dist[v] - 1e-9:
                dist[v] = dist[u] + w
                parent[v] = u

    cycle_vertex = None
    for u, v, w in edges:
        if dist[u] + w < dist[v] - 1e-9:
            parent[v] = u
            cycle_vertex = v
            break

    if cycle_vertex is None:
        return None, 1.0  # No arbitrage exists

    # Step back |V| times to ensure we enter the cycle
    curr = cycle_vertex
    for _ in range(n):
        curr = parent[curr]

    # Reconstruct cycle
    cycle = [curr]
    p = parent[curr]
    while p != curr:
        cycle.append(p)
        p = parent[p]
    cycle.append(curr)
    cycle.reverse()

    multiplier = 1.0
    for i in range(len(cycle) - 1):
        multiplier *= rate_map[(cycle[i], cycle[i+1])]

    return cycle, multiplier
\end{lstlisting}

\subsection*{3.3 Theoretical Complexity Analysis}
\begin{itemize}
    \item \textbf{Time Complexity}: Weight transformation takes $\mathcal{O}(E)$ time. Bellman-Ford performs $|V| - 1$ edge relaxation passes, each scanning all $|E|$ edges, incurring $\mathcal{O}(V \cdot E)$ operations. Cycle extraction takes $\mathcal{O}(V)$. Overall time: $\mathcal{O}(V \cdot E)$. For a complete FX market ($E = V^2$), runtime is $\mathcal{O}(V^3)$.
    \item \textbf{Auxiliary Space Complexity}: $\mathcal{O}(V + E)$ to maintain graph rates, logarithmic weights, distance tables, and parent pointers.
\end{itemize}

\subsection*{3.4 Execution Output \& Verification}
Testing on international currency markets:
\begin{lstlisting}
============================================================
ASSIGNMENT 3: CURRENCY ARBITRAGE DETECTION (BELLMAN-FORD)
============================================================
Scenario 1: Market With Arbitrage Opportunity
  [+] Arbitrage Opportunity Detected!
  Cycle Route:       USD -> EUR -> GBP -> USD
  Product of Rates:  1.022400
  Net Profit Return: +2.24% (Starting with 1,000 USD yields 1022.40 USD)

Scenario 2: Efficient Market (Fair Exchange Rates)
  [-] No Arbitrage Opportunity Detected (Market is efficient and balanced).
Verification Passed: Correctly identified profitable arbitrage loop!
\end{lstlisting}

\newpage

% ==============================================================================
% ASSIGNMENT 4: PROJECT SCHEDULING (DAG)
% ==============================================================================
\section*{Assignment 4: Project Scheduling Using a Directed Acyclic Graph (DAG)}

\subsection*{4.1 Problem Statement and Algorithmic Principle}
\textbf{Problem}: A software engineering project consists of distinct tasks with fixed durations. Precedence constraints dictate that certain tasks cannot begin until prerequisite tasks are fully completed. We seek to determine a valid linear execution order and compute the earliest possible project completion time.

\textbf{Mathematical Formulation}:
Represent project tasks as a Directed Acyclic Graph (DAG) $G = (V, E)$, where vertex $u \in V$ represents a task with duration $d(u) \ge 0$, and directed edge $(u, v) \in E$ indicates that task $u$ must finish before task $v$ can commence.
\begin{enumerate}
    \item \textbf{Topological Ordering}: A linear ordering $\pi = (v_1, v_2, \dots, v_n)$ of $V$ such that for every directed edge $(u, v) \in E$, $u$ appears before $v$ in $\pi$.
    \item \textbf{Earliest Start and Finish Times}:
    \begin{itemize}
        \item Earliest Start Time of task $v$:
        \[
        \text{EST}[v] = \begin{cases}
        0, & \text{if } \text{indegree}(v) = 0 \\
        \max_{(u, v) \in E} \text{EFT}[u], & \text{otherwise}
        \end{cases}
        \]
        \item Earliest Finish Time of task $v$:
        \[
        \text{EFT}[v] = \text{EST}[v] + d(v)
        \]
    \end{itemize}
    \item \textbf{Total Project Completion Time}:
    \[
    T^* = \max_{v \in V} \text{EFT}[v]
    \]
    \item \textbf{Critical Path Analysis}: The critical path is the longest duration chain of dependent tasks. Any delay in a task lying on the critical path directly delays overall project completion time.
\end{enumerate}

\subsection*{4.2 Python Implementation (\texttt{project\_scheduling.py})}
\begin{lstlisting}[language=Python]
from collections import deque
from typing import Dict, List, Tuple

def schedule_project(
    tasks: Dict[str, float],
    dependencies: List[Tuple[str, str]]
) -> Dict:
    adj = {t: [] for t in tasks}
    in_degree = {t: 0 for t in tasks}
    for u, v in dependencies:
        adj[u].append(v)
        in_degree[v] += 1

    # Kahn's Algorithm for Topological Sorting
    queue = deque([t for t in tasks if in_degree[t] == 0])
    topo_order = []
    while queue:
        u = queue.popleft()
        topo_order.append(u)
        for v in adj[u]:
            in_degree[v] -= 1
            if in_degree[v] == 0:
                queue.append(v)

    if len(topo_order) < len(tasks):
        raise ValueError("Cycle detected! Project dependencies must form a DAG.")

    # Calculate Earliest Start Time (EST) and Earliest Finish Time (EFT)
    est = {t: 0.0 for t in tasks}
    eft = {t: float(tasks[t]) for t in tasks}
    predecessors = {t: None for t in tasks}

    for u in topo_order:
        eft[u] = est[u] + tasks[u]
        for v in adj[u]:
            if eft[u] > est[v]:
                est[v] = eft[u]
                predecessors[v] = u

    total_duration = max(eft.values())

    # Trace Critical Path
    last_task = max(tasks.keys(), key=lambda t: eft[t])
    critical_path = []
    curr = last_task
    while curr:
        critical_path.append(curr)
        curr = predecessors[curr]
    critical_path.reverse()

    return {
        'topological_order': topo_order,
        'earliest_start': est,
        'earliest_finish': eft,
        'total_completion_time': total_duration,
        'critical_path': critical_path
    }
\end{lstlisting}

\subsection*{4.3 Theoretical Complexity Analysis}
\begin{itemize}
    \item \textbf{Topological Sorting}: Initializing in-degrees takes $\mathcal{O}(V + E)$. Kahn's queue processing visits every vertex and traverses each edge once, taking $\mathcal{O}(V + E)$ time.
    \item \textbf{EST/EFT Computation}: Iterating through the topological order evaluates each vertex and edge exactly once, costing $\mathcal{O}(V + E)$ time.
    \item \textbf{Total Time Complexity}: $\mathcal{O}(V + E)$ (Linear Time).
    \item \textbf{Auxiliary Space Complexity}: $\mathcal{O}(V + E)$ for adjacency lists, in-degree tables, and schedule arrays.
\end{itemize}

\subsection*{4.4 Execution Output \& Verification}
Testing on a software development task lifecycle:
\begin{lstlisting}
============================================================
ASSIGNMENT 4: PROJECT SCHEDULING USING DAG
============================================================
Valid Execution Order (Topological Sequence):
  Step 1: 1. Requirements Analysis
  Step 2: 2. Architecture Design
  Step 3: 3. Frontend Development
  Step 4: 5. Database Setup
  Step 5: 4. Backend Development
  Step 6: 6. API Integration
  Step 7: 7. Quality Testing (QA)
  Step 8: 8. Production Deployment

Earliest Overall Project Completion Time: 29.0 days
Critical Path (Bottleneck Path):
  1. Requirements Analysis -> 2. Architecture Design -> 5. Database Setup -> 4. Backend Development -> 6. API Integration -> 7. Quality Testing (QA) -> 8. Production Deployment

Task Schedule Details (Days):
  Task                         | Duration  | EST   | EFT  
  ------------------------------------------------------
  1. Requirements Analysis     | 3.0       | 0.0   | 3.0  
  2. Architecture Design       | 4.0       | 3.0   | 7.0  
  3. Frontend Development      | 6.0       | 7.0   | 13.0 
  5. Database Setup            | 3.0       | 7.0   | 10.0 
  4. Backend Development       | 8.0       | 10.0  | 18.0 
  6. API Integration           | 4.0       | 18.0  | 22.0 
  7. Quality Testing (QA)      | 5.0       | 22.0  | 27.0 
  8. Production Deployment     | 2.0       | 27.0  | 29.0 
Verification Passed: Precedence order strictly respected with zero negative slack.
\end{lstlisting}

\newpage

% ==============================================================================
% ASSIGNMENT 5: CAMPUS NAVIGATION (BFS)
% ==============================================================================
\section*{Assignment 5: Campus Navigation Using Breadth-First Search (BFS)}

\subsection*{5.1 Problem Statement and Algorithmic Principle}
\textbf{Problem}: A university campus requires a pedestrian navigation system that computes the path between two buildings having the minimum number of walkway segments. Buildings are represented as vertices and connecting walkways as unweighted edges.

\textbf{Mathematical Formulation \& BFS Optimality}:
Let $G = (V, E)$ be an unweighted undirected graph. The distance $d(u, v)$ between two vertices is defined as the minimum number of edges on any path connecting them.
\begin{itemize}
    \item \textbf{Queue Invariant}: Breadth-First Search (BFS) initializes a FIFO queue with start vertex $s$ at level $0$.
    \item \textbf{Level-by-Level Exploration}: When a vertex at level $k$ is expanded, all its unvisited neighbors are assigned distance $k + 1$ and enqueued.
    \item \textbf{Shortest-Path Guarantee}: Because all edge lengths are equal ($w(e) = 1$), the first time the destination vertex $t$ is popped from the queue, the path traversed contains the absolute minimum number of edges.
\end{itemize}

\subsection*{5.2 Python Implementation (\texttt{campus\_navigation.py})}
\begin{lstlisting}[language=Python]
from collections import deque
from typing import Dict, List, Tuple, Optional

def campus_navigation_bfs(
    walkways: Dict[str, List[str]],
    start_building: str,
    destination_building: str
) -> Tuple[Optional[int], List[str]]:
    if start_building == destination_building:
        return 0, [start_building]

    visited = {start_building}
    parent = {start_building: None}
    queue = deque([start_building])

    while queue:
        curr = queue.popleft()
        if curr == destination_building:
            break

        for neighbor in walkways.get(curr, []):
            if neighbor not in visited:
                visited.add(neighbor)
                parent[neighbor] = curr
                queue.append(neighbor)

    if destination_building not in visited:
        return None, []  # Unreachable

    # Reconstruct path
    path = []
    curr = destination_building
    while curr is not None:
        path.append(curr)
        curr = parent[curr]
    path.reverse()

    return len(path) - 1, path
\end{lstlisting}

\subsection*{5.3 Theoretical Complexity Analysis}
\begin{itemize}
    \item \textbf{Time Complexity}: Every vertex is enqueued and dequeued at most once ($\mathcal{O}(V)$). In doing so, every adjacency list is scanned once ($\mathcal{O}(E)$). Reconstructing the path takes $\mathcal{O}(V)$. Overall time: $\mathcal{O}(V + E)$ (Linear Time).
    \item \textbf{Auxiliary Space Complexity}: $\mathcal{O}(V)$ to store the FIFO queue, visited set, and predecessor reconstruction pointers.
\end{itemize}

\subsection*{5.4 Execution Output \& Verification}
Testing on the SVNIT campus building graph:
\begin{lstlisting}
============================================================
ASSIGNMENT 5: CAMPUS NAVIGATION (BFS)
============================================================
Origin Building:      Main Gate
Destination Building: Central Workshop

Minimum Walkway Segments: 3
Optimal Campus Route:
  Main Gate -> Central Library -> Mechanical Engineering Dept -> Central Workshop
Verification Passed: Path achieves theoretical minimum walkway hop count of 3.
\end{lstlisting}

\newpage

% ==============================================================================
% SYNTHESIS MATRIX & CONCLUSION
% ==============================================================================
\section*{Comprehensive Algorithmic Synthesis \& Comparison Matrix}

Table \ref{tab:comparison_graph} synthesizes the core mathematical properties, algorithmic paradigms, time complexities, auxiliary memory requirements, and optimality guarantees across all five graph optimization problems in Lab Assignment 11.

\begin{table}[H]
\centering
\caption{Synthesis of Algorithmic Paradigms, Complexities, and Optimality Guarantees.}
\label{tab:comparison_graph}
\vspace{0.2cm}
\resizebox{\textwidth}{!}{
\begin{tabular}{lllccc}
\toprule
\textbf{Assignment} & \textbf{Problem Domain} & \textbf{Core Paradigm} & \textbf{Time Complexity} & \textbf{Aux Space} & \textbf{Optimality Guarantee} \\
\midrule
\textbf{Assignment 1} & Delivery Network Design    & Chu-Liu / Edmonds'    & $\mathcal{O}(V \cdot E)$        & $\mathcal{O}(V + E)$ & Exact (Minimum Arborescence) \\
\textbf{Assignment 2} & Emergency Route Finding    & Dijkstra's (Min-Heap) & $\mathcal{O}((V + E) \log V)$   & $\mathcal{O}(V)$     & Exact (Non-negative Weights) \\
\textbf{Assignment 3} & Currency Arbitrage         & Bellman--Ford ($-\ln$) & $\mathcal{O}(V \cdot E)$        & $\mathcal{O}(V + E)$ & Exact (Negative Cycle Bound) \\
\textbf{Assignment 4} & Project Task Scheduling    & Topological Sort/DAG  & $\mathcal{O}(V + E)$            & $\mathcal{O}(V + E)$ & Exact (Longest Path in DAG)  \\
\textbf{Assignment 5} & Campus Walkway Navigation  & Breadth-First Search  & $\mathcal{O}(V + E)$            & $\mathcal{O}(V)$     & Exact (Unweighted Shortest Path) \\
\bottomrule
\end{tabular}
}
\end{table}

\subsection*{Conclusion}
The rigorous computational study of graph optimization algorithms in Lab Assignment 11 highlights several foundational principles of algorithm design:
\begin{enumerate}
    \item \textbf{Directed vs. Undirected Spanning Trees}: While undirected MSTs yield to simple greedy edge ordering (Kruskal/Prim), directed minimum arborescences require recursive cycle contraction (Edmonds' algorithm) because locally optimal incoming edges may form cycles rather than connecting to the root.
    \item \textbf{Edge Weight Signatures Dictate Algorithm Choice}:
    \begin{itemize}
        \item For strictly non-negative weights, Dijkstra's algorithm achieves efficient $\mathcal{O}((V + E) \log V)$ performance by never revisiting settled nodes.
        \item In the presence of negative weights or cycles, Bellman-Ford's $\mathcal{O}(V \cdot E)$ iterative relaxation is essential to detect negative cycles.
    \end{itemize}
    \item \textbf{The Power of Problem Reductions}: The transformation $w(u, v) = -\ln R(u, v)$ maps multiplicative currency arbitrage directly to additive negative cycle detection, demonstrating how mathematical duality bridges financial engineering and graph algorithms.
    \item \textbf{Linear Time Optimality on DAGs and Unweighted Graphs}: By exploiting graph topology (absence of cycles in DAGs via topological sorting, and uniform edge weights via BFS queues), general polynomial-time graph problems collapse to optimal $\mathcal{O}(V + E)$ linear time.
\end{enumerate}

\end{document}

